\documentclass[11pt,a4paper]{article}
\usepackage[margin=2.0cm]{geometry}
\usepackage{amsmath,amssymb,graphicx,booktabs,hyperref,caption,array}
\usepackage{times}
\newcommand{\Hmat}{\mathbf{H}}
\begin{document}

\section*{Supplementary Table S14: Consolidated register of withdrawn, corrected and re-scoped claims}
\label{tab:withdrawals}

This table collects every claim from the earlier versions of this work that has been withdrawn,
replaced, corrected or re-scoped, together with the evidence that motivated the change. It exists
because the retractions are otherwise distributed across the main text, seven Supplementary Notes
and several tables, which makes it hard for a reader to see the full extent of the revision. No
claim listed here is relied on anywhere in the present manuscript.

\vspace{4pt}
\small
\setlength{\tabcolsep}{3pt}
\renewcommand{\arraystretch}{1.3}
\begin{tabular}{@{}p{0.185\textwidth} p{0.125\textwidth} p{0.33\textwidth} p{0.29\textwidth}@{}}
\toprule
\textbf{Original claim} & \textbf{Where it appeared} & \textbf{Why it was withdrawn} & \textbf{Current status / replacement} \\
\midrule

Directional locking in \textbf{21/21} seeds across three tissues, including 15/15 hematopoietic
seeds with $C_{\text{head}}=77.6\%$ and all $64/64$ vanilla directions opposed in every seed.
&
Main text positive control; the opposition-budget result file
&
Recomputation from all 15 surviving paired checkpoints under the published Ridge-readout protocol
gives \textbf{0/15} seeds with $C_{\text{head}}>50\%$ ($C_{\text{head}}$ range $-67\%$ to
$+2.4\%$; opposed directions 4--10 of 64, never 64). The analysis JSON is not derivable from any
surviving checkpoint or preprocessing variant, and its generator script is absent from the
repository and from git history.
&
\textbf{Withdrawn.} Note~5. A fresh 3-seed reproduction locks in 1/3 seeds and only under
early-epoch checkpoint selection; under the unified audit protocol none is below
$\kappa<0.5$. Locking remains plausible but is not demonstrated. \\

Latent effective-rank drop of \textbf{$\sim$30\%} for perturbation-conditioned scVI, read as
evidence that conditioning compresses the latent manifold.
&
Supplementary Figure~S12 (top panel B)
&
The value is a single run. In the paired design---both arms refit at matched seeds---the ratio
averages $1.02\pm0.15$ ($n=4$, $t=+0.2$ against unity), i.e.\ no rank drop.
&
\textbf{Withdrawn.} S12 states the earlier reading explicitly as withdrawn; the family-4 arm
instead contracts slightly ($0.868\pm0.032$, $n=5$). \\

The ``purity paradox'': pure-lineage priors expand the latent while mixed-lineage priors compress it,
separated by $72$ percentage points in the trace form.
&
Main text results section; earlier figures
&
The $72$-pp separation is a matched-scale artefact of the trace protocol. The corrected
predictability form places the pair at $+0.90\%$ against $+0.76\%$ --- a
$0.14$-percentage-point contrast --- and the group-level lineage-purity contrast is
non-significant (pure $n=8$ vs mixed $n=13$, Welch $p=0.74$).
&
\textbf{Withdrawn.} Supplementary Figure~S13; the matched-scale pair, the expansion attributed to
mixed priors, and the reading that prior content tracks the spectral response are all withdrawn. \\

Ridge readout reproduces the analytical Jacobian: dominant-mode agreement
$\cos\theta=0.78\pm0.04$.
&
Methods (Ridge--Jacobian validation)
&
Recomputed from 4 surviving hematopoietic seeds (s9--s12): $\cos\theta=0.45\pm0.26$
(range $0.10$--$0.65$). The originally reported value is not reproducible from the surviving
checkpoints. Ridge- and Jacobian-based $C$ also disagree in sign.
&
\textbf{Replaced} by $\cos\theta=0.45\pm0.26$; the Ridge readout is now described as a proxy for
the decoder's local response map, not the literal weights. Note~1. \\

$\kappa$-based $C$ ranks agree with the Jacobian-based $C$: Spearman $\rho=0.94$.
&
Methods
&
Recomputed on the six G-VAE configurations with a computable analytical Jacobian:
$\rho=-0.03$ ($p=0.96$, $n=6$). Ridge-based $C$ is systematically negative
($-17\%$ to $-584\%$) while Jacobian-based $C$ is positive ($+28\%$ to $+67\%$).
&
\textbf{Withdrawn.} Note~1. Applications needing spectral fidelity are directed to the analytical
Jacobian. \\

GRQ specificity of $10.8$--$12.3\times$ and CTCP specificity of $18.8$--$19.6\times$ as
independent positive evidence for the perturbation framework.
&
Supplementary Table~S6 (earlier version)
&
GRQ values were not reproducible from the surviving checkpoints. Under label-shuffling permutation
tests \emph{with} test-set exclusion, real and shuffled CTCP specificity are statistically
indistinguishable.
&
\textbf{GRQ replaced} by the re-audited per-seed values ($11.3$--$16.8\times$). \textbf{CTCP
re-scoped}: the earlier significance reading was driven by sample-size rebalancing and does not
survive the permutation test. Note~5, Table~S6. \\

Encoder/decoder co-adaptation: encoder weight correlation $r=0.0143$ and decoder correlation
$r=0.964$ reported as co-adaptation evidence.
&
Main text (earlier version)
&
The correlation pair as reported was not reproducible from the surviving checkpoints.
&
\textbf{Replaced} by real recomputed values; \texttt{AUDIT\_NOTES} records the replacement. The
co-adaptation framing is retained only where it survives. \\

The G-VAE encoder is a 3-layer network with widths $512/256/128$.
&
Methods (architecture description)
&
The trained checkpoints define a 2-layer encoder; the width triple did not match the artefacts.
&
\textbf{Corrected} to a 2-layer encoder in \texttt{star\_methods.tex}, verified against the
checkpoint tensors. \\

Decoder-direct perturbation (DDP) reaches $>1{,}000\times$ specificity, presented as a capability.
&
Main text (earlier version); Paper~2
&
The value is a pure mathematical extrapolation rather than a measured readout.
&
\textbf{Re-scoped} to the measured regime: $92$--$156\times$ at
$\|\Delta\mathbf{z}\|=5$--$20$, with the non-physiological limit stated. \\

Neural-crest ($C=56$--$69\%$) and lung-fibrosis ($C=96.7\pm0.3\%$) G-VAE compression figures.
&
Paper~2 figures
&
Produced with the $\kappa$-$C$ penalty term, which is non-differentiable as implemented, so the
figures are not verifiable as spectral results.
&
\textbf{Unverified; withdrawn from use.} Note~5: they should be re-checked with a differentiable
spectral objective before use. \\

The $\kappa$-$C$ regulariser shapes the decoder during training.
&
Methods (earlier framing)
&
\texttt{kappa\_c\_reg()} returns a Python \texttt{float}, so it carries no gradient; a constant
offset cannot move gradients either. Measured $|\Delta\text{grad}|=0$ across three optimisation
steps with and without the term.
&
\textbf{Re-scoped} to a training-objective-agnostic \emph{readout}. The $0/33$ null therefore
compares two differently-initialised models rather than measuring a constrained fit. Note~4,
main text \emph{Scope of the audit}. \\

\bottomrule
\end{tabular}

\vspace{6pt}
\footnotesize\noindent
\textbf{Net effect on the manuscript's claims.} Eleven earlier claims are withdrawn, replaced,
corrected or re-scoped here. The claims the present manuscript rests on --- the band-independent
$0/33$ decoder-gain null, the injection-path finding that prior content does not reach the loss in
the audited implementation, and the diagnostic instrument itself --- are independent of every row
above: none of them depends on a G-VAE positive control, on a Ridge--Jacobian agreement value, or
on a $C$-based ranking.

\end{document}
